\section{Exterior energy for homogeneous dispersion and attraction}
\label{sec:r13-dispersion-structure}

Fix $m>2$, $s=m/2$, and a real even positive elliptic symbol $a$,
$m$-homogeneous and $C^{N+2}$ away from zero:
\[
 c_a|\xi|^m\le a(\xi)\le C_a|\xi|^m.
\]
For $u\in H^s(\mathbb R^N)$ write
\[
 M=\|u\|_2^2,\qquad A=\|a(D)^{1/2}u\|_2^2,\qquad
 P=\operatorname{Im}\int\overline u\nabla u.
\]
We use either the local equation
\begin{equation}\label{r13g:equation}
 iu_t=a(D)u-|u|^{2m/N}u,\qquad
 p=2+2m/N,\qquad E=\frac A2-\frac{\|u\|_p^p}{p},
\end{equation}
or, when $2<m<N$, the Hartree equation
\begin{equation}\label{r13g:hartree-equation}
 iu_t=a(D)u-(K*|u|^2)u,\qquad
 E=\frac A2-\frac{H_K(u)}4,
 \quad H_K(u)=\iint K(x-y)|u(x)|^2|u(y)|^2\,dx\,dy.
\end{equation}
Here $K$ is real, even, $C^1$ away from zero, homogeneous of degree
$-m$, positive definite as a density quadratic form, and
\[
 |K(x)|+|x||\nabla K(x)|\le C_K|x|^{-m}.
\]
The corresponding sharp inequalities and mass excesses are
\begin{align*}
 \|u\|_p^p&\le\frac p2(M/M_c)^{m/N}A,
 &d&=(M/M_c)^{m/N}-1 &&\text{(local)},\\
 H_K(u)&\le2(M/M_c)A,
 &d&=M/M_c-1 &&\text{(Hartree)}.
\end{align*}
The dynamical theorem below uses strongly approximable conservative
$H^s$ solutions. No local theory for arbitrary symbols and dimensions
is silently included in that solution convention.
Positive definiteness here is strict on nonzero real densities in the
interaction space, so its square root is a strictly convex Hilbert
norm. Mere nonnegative semidefiniteness is not the assumed hypothesis.

The exterior kinetic density in both equations is
$|a(D)^{1/2}u|^2$. Convexity of $a$ is unnecessary: all anisotropic
contributions can be confined to an annulus, with an error tending to
zero. We prove that fact directly by Fourier estimates.

\subsection{An exterior kinetic term without convexity}

Put $b=\sqrt a$ and, for $\xi\ne0$, define the bounded degree-zero
symbols
\[
 B_{jk}(\xi)=\frac{\partial_j a(\xi)\xi_k}{a(\xi)}.
\]
Euler's identity gives $\sum_jB_{jj}=m$. These are real even symbols;
no sign is assigned to individual entries.

\begin{lemma}[Cutoff commutations]
\label{lem:r13g-cutoff}
For a fixed smooth compact cutoff $\theta$, $\theta_R(x)=\theta(x/R)$,
and $R\ge1$,
\begin{align}
 \|[B_{jk}(D),\theta_R]b(D)u\|_2
 &\le CR^{-1}\|u\|_{H^s},\label{r13g:zero-order-comm}\\
 \|[b(D),\theta_R]u\|_2
 &\le CR^{-1}\|u\|_{H^s}.\label{r13g:half-comm}
\end{align}
The same estimates hold for constant additions to $\theta_R$.
\end{lemma}
\begin{proof}
Write $h=\xi-\zeta$. The angular derivative bound for $B_{jk}$,
its boundedness, and $s>1$ give
\[
 |B_{jk}(\xi)-B_{jk}(\zeta)|b(\zeta)
 \le C|h||\zeta|^{s-1}.
\]
For $|h|\le|\zeta|/2$ use its homogeneous derivative of degree $-1$;
otherwise use $|\zeta|\le2|h|$. Young's convolution inequality and
$\||h|\widehat\theta_R(h)\|_1=O(R^{-1})$ prove
\eqref{r13g:zero-order-comm}. Similarly,
\[
 |b(\xi)-b(\zeta)|
 \le C|h|\bigl(|\zeta|+|h|\bigr)^{s-1}
\]
gives an upper bound
$CR^{-1}\|u\|_{H^{s-1}}+CR^{-s}\|u\|_2$ for
\eqref{r13g:half-comm}. Constants commute.
\end{proof}

Choose a radial $h_0\in C_c^\infty$ with $0\le h_0\le1$, equal
to one on $B_1$ and zero outside $B_2$. Let
\[
 f_R(x)=x h_0(x/R),\qquad
 L_R=f_R\cdot\nabla+\tfrac12\operatorname{div}f_R,
 \qquad C_R=[a(D),L_R].
\]

\begin{lemma}[Factorization of the principal form]
\label{lem:r13g-principal}
For $g=b(D)u$ and $R\ge1$,
\begin{equation}\label{r13g:factor}
 \langle u,C_Ru\rangle
 =\operatorname{Re}\sum_{j,k}
    \langle g,(\partial_j f_{R,k})B_{jk}(D)g\rangle
       +O(R^{-1}\|u\|_{H^s}^2).
\end{equation}
In particular,
\begin{align}
 \langle u,C_Ru\rangle
 &\le mA(u)-m\int_{|x|>2R}|g|^2
    +C\int_{R/2<|x|<4R}|g|^2
    +CR^{-1}\|u\|_{H^s}^2.
 \label{r13g:exterior-sign}
\end{align}
All constants are unchanged by translation of the cutoffs.
\end{lemma}
\begin{proof}
The Fourier kernel of $C_R$, with $h=\xi-\zeta$, is
\[
 \frac i2\,[a(\xi)-a(\zeta)](\xi+\zeta)_k
                                  \widehat f_{R,k}(h).
\]
The kernel of $\sum_{j,k}(\partial_j f_{R,k})q_{jk}(D)$, where
$q_{jk}(\zeta)=\partial_j a(\zeta)\zeta_k$, is
\[
 i h_j\widehat f_{R,k}(h)\partial_j a(\zeta)\zeta_k.
\]
Their difference is bounded by
\begin{equation}\label{r13g:kernel-remainder}
 C|h|^2(|\zeta|+|h|)^{m-1}|\widehat f_R(h)|.
\end{equation}
This is second-order Taylor expansion: the second derivatives of $a$
extend continuously at zero because $m>2$. To bound the form on $H^s$,
divide the kernel by $\langle\xi\rangle^s\langle\zeta\rangle^s$.
The elementary bound
\[
 \frac{(|\zeta|+|h|)^{m-1}}
      {\langle\xi\rangle^s\langle\zeta\rangle^s}
 \le C\langle h\rangle^{m-1}
\]
follows by splitting $|\zeta|\le2|h|$ and its complement.
Since $\widehat f_R(h)=R^{N+1}\widehat f_1(Rh)$, Schur's test
bounds \eqref{r13g:kernel-remainder} by $CR^{-1}$ as a form
$H^s\times H^s\to\mathbb C$.

Replacing $q_{jk}(D)$ by the $b(D)$ sandwich changes its kernel by
\[
 [b(\zeta)-b(\xi)]b(\zeta)B_{jk}(\zeta)
                      \widehat{\partial_j f_{R,k}}(h).
\]
Its magnitude is at most
$C|h|(|\zeta|+|h|)^{m-1}
|\widehat{\partial_j f_{R,k}}(h)|$.
The same weighted Schur estimate, now using the scale-invariant
coefficient $\partial_j f_{R,k}$, is $O(R^{-1})$.
Taking real parts proves \eqref{r13g:factor}.

Write
\[
 \partial_j f_{R,k}=h_0(x/R)\delta_{jk}+c_{jk,R}.
\]
The coefficients $c_{jk,R}$ are uniformly bounded and supported on
$R<|x|<2R$. The identity $\sum_jB_{jj}=m$ makes the diagonal part
exactly $m\int h_0(x/R)|g|^2$. Choose $\theta_R$ supported on
$R/2<|x|<4R$ and equal to one on the support of every $c_{jk,R}$.
For $c=c_{jk,R}$ and $B=B_{jk}(D)$,
\[
 \langle g,cBg\rangle
 =\langle\theta_Rg,cB(\theta_Rg)\rangle
   +\langle\theta_Rg,c[B,1-\theta_R]g\rangle.
\]
The first term is bounded by $C\|\theta_Rg\|_2^2$, and
Lemma~\ref{lem:r13g-cutoff} bounds the second by
$CR^{-1}\|u\|_{H^s}^2$.
Finally $h_0\le1$ and $h_0=0$ outside $2R$ give
\eqref{r13g:exterior-sign}. No pointwise principal-symbol sign has
been used.
\end{proof}

\subsection{Localized mass and moving virial identities}

\begin{lemma}[Localized mass transport]
\label{lem:r13g-mass-commutator}
For a real smooth compact cutoff $\chi_R(x)=\chi((x-z)/R)$,
$R\ge1$, and $v,w\in H^s$,
\begin{equation}\label{r13g:mass-commutator}
 |\langle v,[a(D),\chi_R]w\rangle|
 \le CR^{-1}\|v\|_{H^s}\|w\|_{H^s}.
\end{equation}
If at a given time $u=e^{i\theta}(X+iY)$, with $X,Y$ real,
$M(Y)\le CdM$, $A(Y)\le CdA(u)$, and the mass is in a fixed
compact positive interval, then for $R A(u)^{1/m}\ge1$,
\begin{equation}\label{r13g:scaled-mass-transport}
 \left|\frac d{dt}\int\chi((x-z)/R)|u|^2\right|
 \le C\sqrt d\,R^{-1}A(u)^{(m-1)/m}.
\end{equation}
The center $z$ in this statement is fixed during differentiation.
\end{lemma}
\begin{proof}
The Fourier kernel of the commutator is
$[a(\xi)-a(\zeta)]\widehat\chi_R(\xi-\zeta)$.
The mean-value bound
\[
 |a(\xi)-a(\zeta)|
 \le C|\xi-\zeta|(|\zeta|+|\xi-\zeta|)^{m-1}
\]
and the weighted Schur estimate in
Lemma~\ref{lem:r13g-principal} give \eqref{r13g:mass-commutator}.
The nonlinear potential in both equations is real and cancels from
the mass derivative. Since $a$ is real even and $\chi_R$ real,
the real-real and imaginary-imaginary terms make no contribution;
only a bilinear expression in $X,Y$ remains. Normalize $A(u)=1$
by critical scaling, apply the commutator bound and
$\|Y\|_{H^s}\le C\sqrt d$, then restore the scale.
The temporal scale contributes $A(u)$, while the cutoff radius becomes
$R A(u)^{1/m}$, giving \eqref{r13g:scaled-mass-transport}.
The chosen constant phase is never differentiated.
\end{proof}

For a locally absolutely continuous center $z(t)$, define
\[
 J_R(t)=2\operatorname{Im}\int f_R(x-z(t))\cdot\overline u\nabla u.
\]
For smooth solutions, direct differentiation gives
\begin{equation}\label{r13g:exact-moving-virial}
 J_R'=2\langle u,C_R^{z(t)}u\rangle+\mathcal N_R
 -2z'(t)\cdot\int Df_R(x-z(t))
                         \operatorname{Im}(\overline u\nabla u).
\end{equation}
The vector field $f_R=xh_0(x/R)$ has symmetric derivative, so this
matrix convention is unambiguous. The nonlinear term is
\[
 \mathcal N_R=-\frac{2(p-2)}p
       \int\operatorname{div}f_R(x-z)|u|^p
       \quad\hbox{for the local power},
\]
and
\[
 \mathcal N_R=\iint|u(x)|^2|u(y)|^2
  [f_R(x-z)-f_R(y-z)]\cdot\nabla K(x-y)\,dx\,dy
       \quad\hbox{for Hartree}.
\]
Their full-dilation values combine with $2mA$ to give $4mE$.
The motion term equals $-2z'\cdot P$ plus an error bounded by
\begin{equation}\label{r13g:motion-defect}
 C|z'|\int_{|x-z|>R}|u||\nabla u|,
\end{equation}
because $Df_R=I$ inside $B_R$. Also, for mass in a fixed compact
interval,
\begin{equation}\label{r13g:virial-endpoint}
 |J_R(t)|\le CR A(u(t))^{1/m}.
\end{equation}
Indeed $\|f_R\|_\infty\le CR$ and interpolation bounds
$\|u\|_2\|\nabla u\|_2$ by $CM^{1-1/m}A^{1/m}$.
All identities are understood in integrated form for the stated
strongly approximable solution class; the Fourier form estimates and
the interaction estimates below justify passage to the limit.

\subsection{Exterior interaction and current estimates}
\label{subsec:r13g-exterior-interaction}

The preceding estimates were stated on a bounded $H^s$ ball. Their
homogeneity gives precise instantaneous kinetic dependence. Fix the mass
in a compact positive interval, put $\lambda=A(u)^{1/m}$, and assume
$R\lambda\ge1$. Rescaling to $A=1$ in
Lemma~\ref{lem:r13g-principal} replaces its remainder by
\begin{equation}\label{r13g:scaled-principal-error}
 CR^{-1}A(u)^{(m-1)/m}.
\end{equation}
For $\zeta_R=0$ on $B_{R/4}$ and $\zeta_R=1$ outside $B_{R/2}$,
Lemma~\ref{lem:r13g-cutoff} similarly gives
\begin{equation}\label{r13g:scaled-cutoff-error}
 \|[a(D)^{1/2},\zeta_R]u\|_2^2
 \le CR^{-2}A(u)^{(m-2)/m}.
\end{equation}
These follow by unitary critical spatial rescaling; the constants do
not require an upper bound on $A(u)$.

Fix a center $z$ and write $g=a(D)^{1/2}u$,
\[
 K_R=\int_{|x-z|>8R}|g|^2,\qquad
 H_R=\int_{R/4<|x-z|<8R}|g|^2,\qquad
 \varepsilon_R=R^{-2}A(u)^{(m-2)/m}.
\]
If $\int_{|x-z|>R/4}|u|^2\le\mu$, then
\begin{equation}\label{r13g:tail-localization}
 M(\zeta_Ru)\le\mu,\qquad
 A(\zeta_Ru)\le C(K_R+H_R+\varepsilon_R).
\end{equation}
Sharp interpolation and interpolation between orders zero and $s$ give
\begin{align}
 \int_{|x-z|>R}|u|^p
 &\le C\mu^{m/N}(K_R+H_R+\varepsilon_R),
 \label{r13g:local-tail-struct}\\
 \int_{|x-z|>R}|u||\nabla u|
 &\le C\mu^{1-1/m}(K_R+H_R+\varepsilon_R)^{1/m}.
 \label{r13g:current-tail-struct}
\end{align}
For the current bound, apply
$\|\nabla w\|_2\le C\|w\|_2^{1-2/m}A(w)^{1/m}$ to
$w=\zeta_Ru$, which equals $u$ on the indicated exterior region.

\begin{lemma}[Hartree virial tail]
\label{lem:r13g-hartree-tail}
In the Hartree setting the difference between the nonlinear localized
virial expression and its full dilation value $-mH_K(u)$ obeys
\begin{equation}\label{r13g:hartree-tail}
 |\text{Hartree defect}|
 \le CR^{-m}M^2+C\mu(K_R+H_R+\varepsilon_R).
\end{equation}
\end{lemma}
\begin{proof}
For $\rho=|u|^2$ the symmetrized nonlinear expression is
\[
 \iint\rho(x)\rho(y)
 [f_R(x-z)-f_R(y-z)]\cdot\nabla K(x-y)\,dx\,dy.
\]
The full dilation vector gives $-mH_K(u)$ by homogeneity.
The difference vanishes if both points lie in $B(z,R)$, and the
uniform Lipschitz bound for $f_R$ bounds its absolute integrand by
$C\rho(x)\rho(y)|x-y|^{-m}$. By symmetry take $|x-z|>R$.
If $|y-z|<R/2$, then $|x-y|\ge R/2$, and this contribution is at
most $CR^{-m}M^2$. Otherwise both points are outside $B(z,R/2)$.
For $w=\zeta_Ru$, HLS and elliptic interpolation give
\[
 \iint|x-y|^{-m}|w(x)|^2|w(y)|^2\,dx\,dy
 \le C\|w\|_{4N/(2N-m)}^4\le CM(w)A(w).
\]
Use \eqref{r13g:tail-localization}. Thus the near-diagonal exterior
interaction sees only the small mass $\mu$, while the separated
core--tail interaction vanishes with $R$.
\end{proof}

\begin{lemma}[Hartree weak energy splitting]
\label{lem:r13g-hartree-splitting}
Suppose $u_n\rightharpoonup v$ in $H^s$ and $u_n\to v$ in
$L^2_{\mathrm{loc}}$. For $w_n=u_n-v$,
\[
 H_K(u_n)=H_K(v)+H_K(w_n)+o(1),\qquad
 E(u_n)=E(v)+E(w_n)+o(1).
\]
\end{lemma}
\begin{proof}
Set $q=4N/(2N-m)$, strictly between $2$ and the $H^s$ Sobolev
exponent. Thus $w_n\to0$ in $L^q_{\mathrm{loc}}$ and is bounded
globally in $L^q$. Approximate $v$ by smooth compact functions to
obtain $\operatorname{Re}(\overline v w_n)\to0$ in $L^{q/2}$.
HLS removes all corresponding terms in the density expansion.
It remains to show
$\iint K(x-y)|v(x)|^2|w_n(y)|^2\,dx\,dy\to0$.
The potential of a smooth compact approximation to $v$ is bounded
because $m<N$, and tends to zero at infinity by the Riesz bound.
Local strong $L^2$ convergence and bounded mass of $w_n$ make this
pairing tend to zero. HLS controls the approximation error.
The kinetic quadratic form splits by weak convergence, proving the
energy identity.
\end{proof}

\subsection{A concrete conservative local flow}

\begin{proposition}[Integer energy regularity in dimensions one to three]
\label{prop:r13g-local-flow}
Let $m=2k$ with integer $k\ge2$, $N=1,2,3$, and let $\mathcal H$
be any separable complex Hilbert space. For every symbol $a$ above,
let $b$ be a common real measurable symbol satisfying
$0\le b(\xi)\le C_b(1+|\xi|^{2k})$. The total-density system
\[
 i\partial_tU=[a(D)+b(D)]U-\|U\|_{\mathcal H}^{4k/N}U
\]
has a unique maximal $C_tH^k_x(\mathcal H)$ solution for every
$H^k(\mathcal H)$ datum. The flow conserves
\begin{gather*}
 M=\int\|U\|_{\mathcal H}^2,\qquad
 P=\operatorname{Im}\int\langle U,\nabla U\rangle_{\mathcal H},\\
 E_b=\frac12\int(a+b)(\xi)\|\widehat U(\xi)\|_{\mathcal H}^2\,d\xi
       -\frac1{\alpha+2}\int\|U\|_{\mathcal H}^{\alpha+2},
 \qquad\alpha=4k/N,
\end{gather*}
and has the usual $H^k$ continuation alternative. The local estimates
are independent of the internal dimension. For the homogeneous case
$b=0$, the localized mass and virial identities also pass by strong
approximation. In particular the scalar dynamical conclusions below
need no additional local-theory assumption throughout this integer-order
class, which includes every positive elliptic even homogeneous
polynomial of degree $2k$. No homogeneous virial identity is asserted
for a general nonzero $b$.
\end{proposition}
\begin{proof}
Write $\alpha=4k/N>k$. On the underlying real Hilbert space, the
radial map $F(Z)=\|Z\|_{\mathcal H}^{\alpha}Z$ is Fr\'echet
$C^{k+1}$. Its derivatives obey, for $1\le j\le k+1$,
\[
 \|D^jF(Z)\|_{\mathrm{op}}
 \le C_{\alpha,j}\|Z\|_{\mathcal H}^{\alpha+1-j},
\]
with constants independent of $\mathcal H$; the derivatives extend
continuously at zero because $\alpha+1-j>0$. These bounds follow
by differentiating the Hilbert norm away from zero: each term is a
product of inner products, vectors, and a power of the norm.
Since $k>N/2$, the Hilbert-valued Sobolev embedding and integer
chain/product estimates, applied to
\[
 F(U)-F(V)=\int_0^1 DF(V+\tau(U-V))(U-V)\,d\tau,
\]
show that $F:H^k(\mathcal H)\to H^k(\mathcal H)$ is locally
Lipschitz. The estimates use operator norms of the displayed
derivatives and pointwise Hilbert norms, so their constants do not
grow with the internal dimension. They extend from finite-component
smooth fields by orthogonal projection and density. The common scalar
group $e^{-it[a(D)+b(D)]}$ is strongly continuous and unitary on
$H^k(\mathcal H)$. Duhamel contraction gives existence, uniqueness,
continuous dependence, and the blowup alternative.

For conservation and approximation, use the self-adjoint Fourier
projection $\Pi_n$ onto $|\xi|\le n$ and solve
\[
 i\partial_tU_n=[a(D)+b(D)]U_n-\Pi_nF(U_n),\qquad
 U_n(0)=\Pi_nU_0,\qquad U_n=\Pi_nU_n.
\]
This projected Hamiltonian flow conserves the restrictions of
$M,P,E_b$: $\Pi_n$ commutes with translations and the linear
dispersion, and its range is invariant under constant phase.
The local Lipschitz bounds are uniform in $n$, so Duhamel and
Gronwall give $U_n\to U$ in $C_tH^k_x(\mathcal H)$ on every compact subinterval
of the lifespan. All three conserved functionals are continuous in
$H^k(\mathcal H)$, using the stated growth bound for $b$, so
conservation passes to the limit. When $b=0$, the same approximation
proves the integrated localized identities: the linear commutators
are continuous $H^k$ forms by the Fourier bounds above, which also
hold componentwise with dimension-independent constants. The nonlinear
projection error vanishes because $F(U_n)\to F(U)$ in $H^k$ and
the fixed-radius first-order virial operator is bounded from $H^k$
to $L^2$. No dispersive decay estimate is required.
\end{proof}

